	\subsection{2. Background and Related
		Work}\label{background-and-related-work}}

\hypertarget{moving-iot-service-composition}{%
	\subsubsection{2.1 Moving IoT Service
		Composition}\label{moving-iot-service-composition}}

Moving IoT services represent a paradigm shift from traditional static
composition: service providers change their spatial positions over time,
creating unique challenges for composition algorithms. The temporal
dimension of service availability requires anticipating future service
positions when making composition decisions {[}1{]}.

A moving crowdsourced service can be modeled as a moving region where
the service provider moves in close proximity to users over time. The
composition problem becomes particularly challenging when considering
spatio-temporal constraints including energy requirements, QoS
parameters, and connectivity ranges. Prior work formalized moving IoT
service composition as a Markov Decision Process where the state
includes service positions, device positions, velocities, and predicted
trajectories {[}1{]}. The action space encompasses service selection,
replacement, addition, and removal operations. The STR-based selection
function provides the fundamental service quality metric.

\hypertarget{actor-critic-deep-reinforcement-learning}{%
	\subsubsection{2.2 Actor-Critic Deep Reinforcement
		Learning}\label{actor-critic-deep-reinforcement-learning}}

Actor-critic algorithms combine value-based and policy-based methods.
The actor learns a stochastic policy directly; the critic estimates the
value function. This architecture reduces variance compared to pure
policy gradient while handling continuous action spaces {[}3{]}.

A2C improves stability through the advantage function, measuring how
much better a particular action is compared to the average. The
advantage is:

\[A(s_t, a_t) = Q(s_t, a_t) - V(s_t) = r_t + \gamma V(s_{t+1}) - V(s_t)\]

Studies on A2C for task scheduling in edge-cloud systems show faster
convergence and better adaptability than DQN {[}3{]}{[}6{]}. Adding LSTM
to A2C handles temporal dependencies in mobility-aware scenarios
{[}3{]}.

\hypertarget{network-architecture-design}{%
	\subsubsection{2.3 Network Architecture
		Design}\label{network-architecture-design}}

Network design affects learning performance. Two architectures dominate:
shared networks where both share feature extraction layers but have
separate output heads, and separate networks with independent parameters
{[}2{]}.

Shared networks reduce parameters, train faster with fewer gradient
computations, and may regularize through shared representation. The risk
is interference between actor and critic updates---gradients from one
affect the other.

Separate networks offer more flexibility for complex states where actor
and critic need different processing. This enables specialized
architectures like LSTM for trajectory encoding in the actor
{[}3{]}{[}9{]}. The trade-off is more computation.

\hypertarget{signal-transmission-reward-str-based-selection-function}{%
	\subsubsection{2.4 Signal Transmission Reward (STR) Based Selection
		Function}\label{signal-transmission-reward-str-based-selection-function}}

The service selection uses a hierarchical model where capacity derives
from the Signal Transmission Reward (STR), which depends on Euclidean
distance between the consumer and service provider.

\textbf{Step 1 - Euclidean Distance Calculation}: The distance between
service \(i\) and user \(j\) is:
\[d_{ij} = \sqrt{(x_i - x_j)^2 + (y_i - y_j)^2}\]

\textbf{Step 2 - STR Calculation (Exponential Attenuation)}: The STR
models signal transmission probability with distance-based exponential
attenuation:
\[STR(d_{ij}) = \begin{cases} 1 & \text{if } d_{ij} \leq R_c \\ e^{-k \cdot (d_{ij} - R_c)} & \text{if } d_{ij} > R_c \end{cases}\]

Where: - \(d_{ij}\) is the Euclidean distance between service \(i\) and
user \(j\) - \(R_c\) is the confident radius defining the region where
full transmission is guaranteed (200-300 m) - \(k\) is the decay factor
determining the rate of signal attenuation (0.01-0.05)

\textbf{Step 3 - Capacity Calculation}: Using the Shannon-Hartley
theorem, the channel capacity depends on STR:
\[C_{ij} = B \cdot \log_2(1 + STR(d_{ij}) \cdot SNR_{max})\]

Where: - \(B\) is the bandwidth in Hz (10 MHz) - \(SNR_{max}\) is the
maximum SNR at zero distance (1000 or 30 dB)

\textbf{Step 4 - Reward Calculation (Capacity Based)}: The reward comes
from capacity:
\[R(s_t, a_t) = \begin{cases} +C_{total} & \text{if } C_{total} > C_{min} \\ -1 & \text{if } C_{total} \leq C_{min} \end{cases}\]

Where \(C_{total} = \sum_{i \in C} C_{ij}\) is the total capacity of the
composition and \(C_{min}\) is the minimum required capacity.

The overall service selection function combines STR-derived capacity
with service attributes:
\[i^* = \arg\max_{i \in S} \left[ C_{ij} \cdot E_i \cdot T_{available}^i \right]\]

This hierarchical model ensures that proximity (lower distance leads to
higher STR, which translates to higher capacity) drives the composition
decision while also considering energy and temporal factors.

\hypertarget{recent-advances-in-drl-for-service-composition}{%
	\subsubsection{2.5 Recent Advances in DRL for Service
		Composition}\label{recent-advances-in-drl-for-service-composition}}

Research applies DRL to service composition in various ways. Multi-user
edge service orchestration uses DRL for QoS optimization through
parametric combinatorial action modeling {[}11{]}. Graph RL enables
dependency-aware microservice deployment in edge environments, using
graph convolutional networks for complex call graphs {[}12{]}.

Multi-agent systems with DRL scale IoT service composition. The DRL-MAS
framework combines decentralized multi-agent decision-making for
scalability, energy efficiency, and responsiveness {[}13{]}. Graph
convolutional multi-agent DRL enables dynamic service function chain
deployment with multi-objective optimization across delay and resource
utilization {[}15{]}.

Aerial-terrestrial network integration uses DRL for service composition
with trajectory prediction {[}15{]}. Collective DRL enables intelligent
sharing across edge nodes using soft actor-critic {[}16{]}.

\hypertarget{surveys-on-iot-service-composition}{%
	\subsubsection{2.6 Surveys on IoT Service
		Composition}\label{surveys-on-iot-service-composition}}

Two surveys analyze IoT service composition. Asghari et al.~{[}26{]}
reviewed literature from 2012-2017. Hamzei et al.~{[}27{]} surveyed
approaches as framework, service-oriented architecture/RESTful,
heuristic, and model-based. Key challenges: scalability (45.4\% of
articles), execution time (36.3\%), cost (27.2\%), reliability (22.7\%).
Arellanes et al.~{[}28{]} evaluated scalability---dataflow,
orchestration, and choreography do not fully satisfy scalability; DX-MAN
shows promise.

\hypertarget{theoretical-foundations}{%
	\subsubsection{2.7 Theoretical
		Foundations}\label{theoretical-foundations}}

Actor-critic convergence has been studied. Finite-time analysis shows
single-timescale actor-critic with linear function approximation finds
an \(\epsilon\)-approximate stationary point with
\(\mathcal{O}(\tilde{\epsilon}^{-2})\) sample complexity {[}22{]}. This
supports applying A2C to moving service composition.

MLMC-NAC achieves \(\tilde{\mathcal{O}}(1/\sqrt{T})\) convergence for
average-reward MDPs without requiring mixing and hitting times {[}21{]}.

For multi-objective RL, MOAC provides finite-time convergence and sample
complexity independent of objective count {[}22{]}. With our
multi-component reward, this assures convergence despite complex
objectives.

Single-loop actor-critic with compatible function approximation achieves
optimal sample complexity by eliminating critic approximation error
{[}24{]}. This fits our online service composition.

Proactive composition anticipates future states rather than just
reacting. Latency-aware and proactive service placement uses exponential
smoothing for QoS prediction in mobile edge {[}17{]}. Spatial-temporal
neural networks for connected vehicles achieve 6\% higher prediction
accuracy and 10\% lower service dropping through gated recurrent units
and graph convolutional layers {[}4{]}. ESPD-LP improves data
transmission by 41\% through bidirectional matching across MEC servers
{[}19{]}.

\hypertarget{related-work}{%
	\subsubsection{2.8 Related Work}\label{related-work}}

Service composition has been extensively studied in static environments.
Traditional approaches use QoS-based optimization, selecting services
that maximize composite QoS while satisfying constraints {[}11{]}. These
methods assume fixed service locations and ignore mobility.

\textbf{Moving IoT Service Composition}: Recent work addresses the
unique challenges of moving services. The spatio-temporal nature means
services are available only when they overlap with the user in both
space and time. Fixed services (like WiFi at a coffee shop) remain at
known locations, while moving services (like a smartphone hotspot while
walking) change position continuously. Prior work formalized this as
trajectory-based composition where both provider and consumer have GPS
trajectories {[}1{]}. The challenge is that co-movement patterns must be
discovered---finding services that overlap with the user's path at each
timestep.

\textbf{Deep Reinforcement Learning for Service Composition}: DRL has
emerged as a powerful approach for service composition in dynamic
environments. Unlike traditional optimization, DRL learns a policy
through interaction with the environment, adapting to changing
conditions. DQN-based approaches have been applied to mobile service
composition {[}1{]}, but suffer from overestimation bias and reactive
decision-making. Policy gradient methods like A2C address these
limitations by learning the policy directly.

\textbf{Actor-Critic Methods for Edge Computing}: A2C has shown promise
in edge computing scenarios. Studies on A2C for task scheduling in
edge-cloud systems demonstrate faster convergence and better
adaptability than DQN {[}3{]}{[}6{]}. Adding LSTM to A2C handles
temporal dependencies in mobility-aware scenarios {[}3{]}, enabling the
agent to learn from trajectory patterns.

\textbf{Network Architecture Design}: The choice between shared and
separate networks impacts performance. Shared networks reduce parameters
and train faster with fewer gradient computations. Separate networks
offer flexibility for specialized processing like LSTM trajectory
encoding {[}3{]}{[}9{]}. Our work compares both architectures for moving
IoT service composition.

\hypertarget{literature-review}{%
	\subsubsection{2.9 Literature Review}\label{literature-review}}

This section reviews relevant literature across four areas: moving IoT
service composition, deep reinforcement learning for services,
actor-critic methods for edge computing, and trajectory-based service
discovery.

\textbf{Moving IoT Service Composition}: Traditional service composition
assumes static services at known locations {[}11{]}. Moving IoT services
fundamentally change this assumption---services change position
continuously, requiring composition decisions at each timestep. Prior
work formalized moving IoT service composition as a trajectory-based
problem where both provider and consumer have GPS trajectories {[}1{]}.
Key challenges include: (1) discovering services that overlap with the
user in both space and time, (2) ensuring service continuity when
trajectories only partially overlap, and (3) scaling to large trajectory
datasets without expensive indexing.

The Signal Transmission Reward (STR) model provides a hierarchical
approach: distance determines STR, STR determines capacity, and capacity
determines service quality {[}1{]}. This model captures the physical
reality that wireless signal strength degrades with distance.

\textbf{Deep Reinforcement Learning for Service Composition}: DRL has
emerged as the dominant approach for dynamic service composition.
DQN-based methods {[}1{]} learn Q-values for service selection but
suffer from overestimation bias---the Q-values systematically
overestimate true action values, leading to suboptimal selection when
valid services are scarce {[}7{]}. Double DQN addresses this by
separating action selection from evaluation but still operates
reactively.

Policy gradient methods including REINFORCE and Actor-Critic address
DQN's limitations by learning the policy directly. Policy gradient
methods are unbiased and naturally handle discrete action spaces.
However, pure policy gradient suffers from high variance. Actor-critic
methods reduce variance by combining policy gradient with a value
function baseline.

\textbf{Actor-Critic Methods for Edge Computing}: A2C has demonstrated
success in edge computing scenarios. Chen et al.~{[}3{]} applied
LSTM-based A2C to network slicing with user mobility, showing faster
convergence than DQN. Liu et al.~{[}6{]} proposed A2C-DRL for dynamic
scheduling in edge-cloud environments. These works confirm A2C's
advantages in mobility-aware scenarios.

Theoretical analysis provides convergence guarantees. Zhang et
al.~{[}21{]} showed \(\tilde{\mathcal{O}}(1/\sqrt{T})\) convergence for
average-reward MDPs. Xiao et al.~{[}22{]} provided finite-time analysis
for multi-objective actor-critic. These results assure convergence
despite complex reward structures.

\textbf{Trajectory-Based Service Discovery}: Co-movement pattern
discovery identifies services that overlap with user trajectories.
Traditional methods use flock, convoy, swarm, and group patterns
{[}16{]}{[}21-24{]}. However, these centralized index-based approaches
degrade as datasets scale {[}20{]}.

Parallel flock-based approaches using MapReduce address scalability
{[}1{]}. The spatio-temporal MapReduce first prunes trajectories
temporally, then filters spatially to find candidate services. Our work
builds on this by using DRL to learn composition without indexing.

\begin{center}\rule{0.5\linewidth}{0.5pt}\end{center}

\hypertarget{proposed-a2c-based-framework}{%
